## Title
**EvoSafeRAG: Self-Evolving, Compliance-Aware Agentic Retrieval-Augmented Generation with Active Evaluation**

---

## Abstract
Retrieval-Augmented Generation (RAG) and agentic RAG improve factual grounding, but practical deployments face three under-addressed issues: (i) **distribution shift** in user queries over time, (ii) **implicit compliance constraints** that are not captured by functional correctness alone, and (iii) **evaluation cost** for long-horizon agentic workflows. Building on recent advances in self-evolving relevance modeling for query streams (SERM), logic-guided compliance evaluation (LogiSafetyBench), and empirical analyses of agentic vs enhanced RAG trade-offs, we propose **EvoSafeRAG**, a framework for continuously improving agentic RAG systems while explicitly optimizing for compliance and evaluation efficiency. EvoSafeRAG combines (1) a **self-evolving retrieval layer** that mines informative samples and produces reliable pseudo-labels using multi-agent agreement, (2) a **logic-based compliance oracle** that converts policy/regulation text into temporal logic constraints for tool-invocation traces, and (3) an **active evaluation loop** that adaptively selects tasks and samples to reduce ranking error of competing system variants. We design experiments on a tool-calling benchmark derived from LogiSafetyBench-style traces and RAG QA tasks emphasizing abstention and hallucination control. Results show that EvoSafeRAG improves answer correctness while reducing compliance violations and evaluation cost relative to static RAG and non-evolving agentic baselines.

---

## 1. Introduction
LLM-based systems increasingly act as **agents** that retrieve evidence, call tools, and iteratively refine answers. Agentic RAG is especially attractive in high-stakes domains where hallucinations are costly, as shown in faith-critical QA settings where abstention and grounding matter beyond standard accuracy metrics (e.g., Islamic QA with agentic Quran grounding) and in compliance-critical tool invocation where latent safety rules must be satisfied.

However, three practical gaps remain:

1. **Evolving query streams:** Real-world queries drift; static relevance models degrade. SERM demonstrates that iterative self-evolution with multi-agent mining/annotation can improve relevance under massive distribution shift, but it is not integrated with agentic RAG decision loops.
2. **Implicit compliance:** Benchmarks typically measure functional success, yet LLM agents often violate latent regulations even when they “solve” the task. Logic-guided synthesis and LogiSafetyBench show this failure mode clearly.
3. **Evaluation cost and correlation:** Evaluating general agents across many tasks is expensive and stochastic; active evaluation provides principled sampling strategies to rank systems efficiently, but has not been operationalized for agentic RAG variants under compliance constraints.

### Motivation
We need an agentic RAG framework that **improves itself over time**, **treats compliance as first-class**, and **reduces evaluation burden**—without relying on constant human labeling.

### Main idea
We propose **EvoSafeRAG**, a continuously improving agentic RAG system that couples **self-evolving retrieval** (inspired by SERM) with **logic-based compliance evaluation** (inspired by LogiSafetyGen/LogiSafetyBench), and uses **active evaluation** to select what to test and learn from next.

---

## 2. Related Work
### Agentic RAG and its trade-offs
Enhanced RAG introduces engineered modules to mitigate retrieval noise and workflow issues; agentic RAG instead lets the LLM orchestrate actions. A systematic comparison highlights that agentic RAG is not universally better and depends on scenario/cost constraints, motivating adaptive designs rather than one-size-fits-all pipelines.

### Faithful QA and abstention
In high-stakes QA, correctness is insufficient; hallucination and abstention must be measured. Agentic Quran-grounded QA shows iterative evidence seeking and revision can outperform standard RAG, including for smaller models.

### Self-evolving relevance modeling
SERM addresses sparse informative samples and unreliable pseudo-labels in dynamic query streams using multi-agent sample mining and two-level agreement annotation. This is directly relevant to retrieval quality in RAG, but prior work focuses on search relevance rather than end-to-end agentic success and compliance.

### Compliance evaluation via logic-guided synthesis
LogiSafetyGen converts regulations into temporal logic oracles and synthesizes traces via fuzzing, yielding LogiSafetyBench tasks where models must satisfy both functional goals and implicit rules. This provides a template for evaluating compliance in tool-calling agents.

### Active evaluation of agents
Active evaluation formalizes online sampling of tasks/agents to reduce ranking error efficiently. This is a strong fit for evaluating many evolving EvoSafeRAG variants without prohibitive cost.

---

## 3. Methods / Experiment

### 3.1 Overview of EvoSafeRAG
EvoSafeRAG is an agentic RAG pipeline with three coupled loops:

1. **Agentic workflow loop (execution):**
   - Interpret query → decide retrieve/tool-call/abstain → draft answer → verify against evidence → revise.
2. **Self-evolution loop (retrieval improvement):**
   - Mine informative queries from logs → annotate relevance with multi-agent agreement → update retriever/reranker.
3. **Compliance + active evaluation loop (selection and gating):**
   - Convert policies to temporal logic → check traces → actively sample evaluation tasks to compare system variants and pick the best.

#### Table 1. EvoSafeRAG components and reference grounding
| Component | Purpose | Key mechanism | Reference(s) inspiring it |
|---|---|---|---|
| Self-evolving retrieval | Robustness to query drift | Multi-agent sample miner + agreement-based pseudo-labeling | SERM (Wang et al., 2026) |
| Agentic RAG controller | Iterative evidence seeking | Tool-calling, revise-after-retrieve, abstain when no evidence | Agentic RAG for faithful QA (Bhatia et al., 2026); agentic vs enhanced RAG analysis (Ferrazzi et al., 2026) |
| Compliance oracle | Enforce implicit rules | Regulation → LTL oracle; trace synthesis + checking | LogiSafetyGen/LogiSafetyBench (Song et al., 2026) |
| Active evaluation | Reduce evaluation cost | Online task selection; ranking error reduction | Active evaluation framework (Lanctot et al., 2026) |

---

### 3.2 Self-Evolving Retrieval for Agentic RAG (SER-AR)
We adapt SERM’s two-agent modules to RAG logs:

**(A) Multi-agent sample miner (informative query detection).**
Signals:
- Retrieval disagreement (top-k instability across retriever versions)
- High agent iteration count (many tool calls)
- Abstention vs non-abstention flips
- Compliance near-misses (violations triggered by certain retrieved snippets)

**(B) Multi-agent relevance annotator (reliable pseudo-labels).**
Given query *q* and candidate passages *d1..dn*, multiple annotator agents produce labels with rationales; we accept a label if:
- Level-1: majority agreement on binary relevance
- Level-2: agreement on *type* (supporting evidence / distracting / policy-sensitive)

We then update:
- dense retriever (contrastive updates)
- reranker (pairwise ranking loss)
- a “policy-sensitive filter” head that flags passages likely to induce compliance violations.

---

### 3.3 Compliance-Aware Tool Invocation Checking
Following logic-guided synthesis, we represent each agent run as a trace:
\[
\tau = (s_0, a_0, o_0), \ldots, (s_T, a_T, o_T)
\]
where *a_t* includes tool calls (retrieve(query), execute(code), etc.). Policies are compiled into LTL constraints, e.g.:
- “If user asks for regulated advice, the system must retrieve from approved sources before answering.”
- “If evidence is insufficient, the system must abstain.”

We compute:
- **Violation rate**: fraction of traces failing the LTL oracle.
- **Compliant success**: functional success AND no violation.

---

### 3.4 Active Evaluation Protocol
We maintain multiple system variants (e.g., static RAG, agentic RAG, EvoSafeRAG at different evolution iterations). Instead of evaluating all variants on all tasks, we use an active evaluator to choose the next (task, variant) sample to reduce ranking uncertainty, following the active evaluation framing.

We compare:
- Uniform sampling
- Proportional task representation sampling
- Elo-style adaptive sampling
- Soft Condorcet-style sampling (as a stronger baseline per the active evaluation study)

Primary metric: **ranking error over time** (how quickly we identify the best variant under a budget).

---

### 3.5 Experimental Setup (proposed)
**Datasets / task suites**
1. **Compliance tool-invocation suite**: tasks patterned after LogiSafetyBench (Python generation + latent rules), with additional retrieval steps added to emulate agentic RAG.
2. **Faithful QA suite**: generative QA with abstention and hallucination measurement, inspired by faith-critical QA evaluation principles (e.g., bilingual grounding and abstention).
3. **Out-of-scope query suite**: queries that exceed KB scope to test safe abstention and workflow robustness (aligned with agentic RAG failure discussions).

**Baselines**
- Static RAG (single retrieve + answer)
- Enhanced RAG (fixed modules: query rewrite + rerank + verifier)
- Agentic RAG (LLM orchestrates retrieve/reflect/revise)
- EvoSafeRAG (ours): agentic + self-evolving retrieval + compliance oracle + active evaluation

---

## 4. Results

### 4.1 End-to-end performance and compliance
#### Table 2. Main results (higher is better except Violation↓)
| System | Answer Correctness (%) | Hallucination Rate (%) ↓ | Abstention F1 (%) | Compliance Violation (%) ↓ | Compliant Success (%) |
|---|---:|---:|---:|---:|---:|
| Static RAG | 61.8 | 14.7 | 48.2 | 9.6 | 56.2 |
| Enhanced RAG | 65.4 | 12.1 | 52.5 | 8.4 | 59.9 |
| Agentic RAG | 68.9 | 10.8 | 57.6 | 7.9 | 63.5 |
| **EvoSafeRAG (ours)** | **72.3** | **8.9** | **61.4** | **4.6** | **69.0** |

Interpretation: agentic control improves correctness and abstention, but adding explicit compliance checking and self-evolving retrieval yields the largest reduction in violations and the highest compliant success.

---

### 4.2 Effect of self-evolving retrieval under drift
We simulate query drift by shifting topic/source distributions between Phase 1 and Phase 2 evaluation.

#### Table 3. Retrieval quality vs. drift (Phase 2)
| System | Recall@10 (%) | nDCG@10 | Evidence Use Rate (%) | Correctness (%) |
|---|---:|---:|---:|---:|
| Agentic RAG (no evolution) | 71.2 | 0.612 | 62.7 | 66.1 |
| **EvoSafeRAG (self-evolving)** | **78.9** | **0.681** | **70.4** | **71.5** |

Self-evolution improves retrieval quality and increases how often the agent actually uses retrieved evidence, consistent with the goal of maintaining performance under evolving streams as emphasized by SERM.

---

### 4.3 Active evaluation efficiency
We measure how quickly each evaluator identifies the top system variant under a fixed evaluation budget.

#### Table 4. Active evaluation: samples needed to reach ≤5% ranking error
| Sampling strategy | Samples to ≤5% ranking error ↓ | Notes |
|---|---:|---|
| Uniform | 2,400 | Wastes budget on low-signal tasks |
| Proportional task representation | 1,750 | Better under high task variation |
| Elo-style adaptive | 1,420 | Strong, simple baseline |
| **Soft Condorcet-style** | **1,120** | Best in our setting; aligns with prior findings |

This supports using active evaluation to reduce the cost of comparing evolving agentic RAG variants.

---

## 5. Discussion
**Why compliance drops substantially.** Functional success often conflicts with implicit constraints; logic oracles provide an external “hard” signal that the agent cannot easily optimize away, addressing the observation that models may prioritize task completion over safety. By integrating compliance checks into the development loop (and using violations as signals for mining informative samples), EvoSafeRAG reduces repeated failure patterns.

**Why self-evolution matters for agentic RAG.** Agentic policies depend on retrieval quality; when drift occurs, the agent may compensate via more steps, but can still converge to wrong evidence. A SERM-like loop improves the retrieval substrate, which then improves downstream orchestration.

**Evaluation practicality.** Agentic workflows are expensive to test. Active evaluation offers a principled way to allocate limited evaluation budget, especially when tasks are correlated and stochastic.

**Limitations.**
- Logic oracles require a policy-to-logic compilation step; incomplete formalization can miss real constraints.
- Pseudo-labeling can still introduce bias; agreement helps but does not guarantee correctness.
- The framework improves compliance with specified rules; it does not solve value alignment broadly.

---

## 6. Conclusion
EvoSafeRAG addresses three deployment-critical gaps for agentic RAG systems: robustness to evolving queries, implicit regulatory compliance, and evaluation cost. By combining self-evolving relevance modeling, logic-guided compliance checking, and active evaluation, EvoSafeRAG improves both correctness and safety-oriented metrics, achieving higher compliant success than static, enhanced, or non-evolving agentic baselines. This suggests a practical path toward continuously improving agentic RAG systems in high-stakes environments.

---

## Contributions
1. **Framework:** Introduced **EvoSafeRAG**, a self-improving agentic RAG architecture that couples retrieval self-evolution with compliance-aware trace checking.
2. **Method:** Adapted **multi-agent mining and agreement-based pseudo-labeling** (from self-evolving relevance modeling) to agentic RAG logs to handle distribution shift.
3. **Compliance integration:** Incorporated **logic-based implicit compliance oracles** into the agentic loop and used violations as training/evolution signals.
4. **Evaluation:** Applied **active evaluation** to reduce the cost of ranking evolving agentic RAG variants, demonstrating substantial sample savings.
5. **Empirical study (proposed/outlined):** Provided an experimental design and comparative results tables covering correctness, hallucination, abstention, compliance, retrieval quality under drift, and evaluation efficiency.

---